% Generated draft: requires mathematical and source review.
\documentclass[11pt]{article}
\usepackage[margin=25mm]{geometry}
\usepackage{fontspec}
\setmainfont{Latin Modern Roman}
\usepackage{amsmath,amssymb,mathtools,unicode-math}
\setmathfont{Latin Modern Math}
\usepackage{enumitem,needspace,xurl,hyperref}
\setlength{\parindent}{0pt}
\setlength{\parskip}{5pt}
\setlength{\emergencystretch}{3em}
\newcommand{\N}{\mathbb{N}}
\newcommand{\Z}{\mathbb{Z}}
\newcommand{\Q}{\mathbb{Q}}
\newcommand{\R}{\mathbb{R}}
\newcommand{\C}{\mathbb{C}}
\newcommand{\F}{\mathbb{F}}
\newcommand{\A}{\mathbb{A}}
\newcommand{\PP}{\mathbb{P}}
\newcommand{\E}{\mathbb{E}}
\newcommand{\eps}{\varepsilon}
\newcommand{\dd}{\,\mathrm{d}}
\begin{document}
\section*{UnsolvedMath ID: 20000822}
\textbf{Source problem code:} catalogue entry
\subsection{Definitions and mathematical statement}
For a polynomial ring $S$ with an abelian-group grading and an admissible Hilbert function $h$, the multigraded Hilbert scheme parametrizes homogeneous ideals whose quotient graded pieces have the prescribed dimensions, with the corresponding flat-family condition. A fat point here means $\operatorname{Spec}A$ for a finite-dimensional local algebra $A$ with nonzero nilradical, hence length greater than one and a single reduced support point.

\textbf{Problem 21.} Does some multigraded Hilbert scheme have a connected component isomorphic, as a scheme, to a fat point? The source does not fix the grading, Hilbert function, ring, field or characteristic. Merely having a nonreduced component of positive dimension does not answer the question.

\textbf{Reported fixed-locus criterion.} A grading refinement can realize a refined multigraded Hilbert scheme inside the scheme-theoretic fixed locus of the additional grading torus. If the fixed local scheme at an isolated point is Artin local and nonreduced, and lies in the corresponding refined-Hilbert-function locus, it gives the requested fat connected component. Isolation and the fixed scheme structure are both required.

For a torus-graded affine ring $R=\bigoplus_\chi R_\chi$, the fixed subscheme uses the quotient
\[
 R\big/\langle R_\chi:\chi\ne0\rangle,
\]
not the invariant subring $R^T$, which instead defines an affine quotient. The report applies this distinction to a particular nonreduced toric Hilbert chart.

\textbf{Reported chart obstruction.} Over $\C$ let
\[
 R=\C[z_5,z_{10},z_{11},z_{13}]
 \big/\big((z_5z_{10}^2-z_{13})z_{10}z_{11}^2\big),
\]
and consider the origin, corresponding to the designated monomial point. The $z_{11}^2$ factor gives the chart's doubled irreducible component. For an additional grading torus diagonal in these minimal coordinates, with character weights $w_i$, preservation of the displayed relation imposes
\[
 w_{13}=w_5+2w_{10}.
\]
The scheme-theoretic fixed quotient is obtained by setting to zero every coordinate of nonzero character. The aid proves that every such compatible diagonal action has fixed locus through the origin either a reduced zero-dimensional point or a positive-dimensional scheme. It is never an Artin nonreduced local scheme there.

Thus refining the grading diagonally in these four displayed coordinates cannot isolate the doubled component as a fat point. The conclusion is limited to that chart and those diagonal refinements; it does not exclude other grading actions, other charts or other multigraded Hilbert schemes. The aid reports a failed candidate construction and a criterion, not an answer to the general existence problem.
\subsection{Short English statement}
The source seeks a connected component that is entirely a nonreduced point. The aid gives a scheme-theoretic fixed-locus route, then shows that every compatible diagonal refinement of a standard four-parameter toric Hilbert chart fails: its fixed locus is either a reduced point or remains positive-dimensional. Other constructions remain unaddressed.
\subsection{Sources}
\begin{itemize}
\item Local UnsolvedMath registry, permanent ID $20000822$, \texttt{catalogue entry}; canonical \texttt{aim-arithmetic-geometry-notes.json}, \texttt{notes[67]}.
\item AIM workshop ``Components of Hilbert schemes,'' PDF Problem 21: \url{https://aimath.org/WWN/hilbertschemes/hilbertschemes.pdf}.
\item Stillman--Sturmfels--Thomas chart, as identified in the supplied summary by its explicit four-variable ring. No full title, location or bibliographic reference is supplied.
\item Research attempt 1, \texttt{reduction}, run status \texttt{open}, internal rigor \texttt{proved}, low novelty confidence; current registry partially solved. Explicitly a machine-generated research aid, not peer reviewed; independently verify before citation. Original literature field empty.
\end{itemize}
\end{document}
